\section{NN-S016: an explicit Chow basis and sphere homotopy type with a counting-hard rank}

\paragraph{Statement and exact inputs.}
Let $X_\Phi/\mathbb Z$ be the smooth sparse cubic of NN-C027, and write
\begin{gather*}
 X_\Phi:\quad t^2-t=\sum_{i=1}^h y_i f_i(x_1,\ldots,x_v),\\
 Y_\Phi=V(f_1,\ldots,f_h)=\coprod_{a\in\mathcal S}\operatorname{Spec}\mathbb Z,
 \qquad |\mathcal S|=N=\#\operatorname{SAT}(\Phi).
\end{gather*}
Here $v,h\ge1$, the retained equation is $f_1=x_1^2-x_1$, and
$\mathcal S$ consists of the uniquely extended satisfying Boolean assignments.
For every field $k$, the variety $X=X_\Phi\times k$ is $k$-rational,
smooth and geometrically integral. For $a\in\mathcal S$ and $e\in\{0,1\}$,
put
\[
 D_{a,e}=V(x-a,t-e)\subset X,\qquad D_{a,e}\simeq\mathbb A^h_k.
\]
These are closed smooth subvarieties of codimension $v$. The full integral
Chow ring is
\begin{equation}\label{eq:S016-chow}
 \operatorname{CH}^*(X)\simeq
 \mathbb Z[u_a:a\in\mathcal S]/(u_a u_b:a,b\in\mathcal S),
 \qquad \deg u_a=v,
\end{equation}
where $u_a=[D_{a,0}]$ and $[D_{a,1}]=-u_a$.
In particular there is no integral torsion, and the only nonzero positive
codimension Chow group is $\operatorname{CH}^v(X)=\mathbb Z^N$.
Every element of this group has an effective cycle representative.
For every $\ell\ne\operatorname{char}k$, the cycle-class map is an isomorphism
\begin{equation}\label{eq:S016-cycle}
 \operatorname{CH}^j(X)\otimes\mathbb Q_\ell
       \xrightarrow{\ \sim\ }
 H^{2j}(X_{\bar k},\mathbb Q_\ell(j))
 \quad\hbox{for all }j\ge0.
\end{equation}
All odd geometric cohomology vanishes. The positive codimension generators
are already defined over the ground field.

Over $\mathbb C$ there is also a homotopy equivalence
\begin{equation}\label{eq:S016-wedge}
                  X_\Phi(\mathbb C)\simeq\bigvee_{a\in\mathcal S}S^{2v},
\end{equation}
where the empty wedge is a point. Thus exact computation of the positive
Chow rank, the nonzero reduced Betti number, or the number of spheres is
\#P-complete on this encoded family. Noncontractibility is NP-complete and
contractibility is coNP-complete on the same family. In addition, $N=0$ if and only if $X_\Phi$ is isomorphic to
$\mathbb A^{v+h}_{\mathbb Z}$, and if and only if its fiber over any
specified field is isomorphic to affine space. Affine-space recognition
is therefore coNP-complete on this family. These assertions concern
the supplied circuit-and-equation encoding; no such upper complexity bound
for arbitrary affine cubics is claimed.

Besides NN-C027, the algebraic inputs are localization for Chow groups
(Stacks Project, Lemma~42.19.3, Tag~02RX), affine-bundle surjectivity
(Lemma~42.32.1, Section Tag~02TS), and homotopy invariance for vector bundles
(Lemma~42.36.3, Tag~02TY). The cohomological inputs are SGA~4, Expos\'e~XVI,
Theorem~3.7 and Remark~3.10(a), and Expos\'e~XVIII, Theorem~3.2.5 and
equation~(3.2.6.2); SGA~5, Expos\'e~VII, Proposition~1.1 and Corollary~1.2
give affine-space invariance. The cycle normalization is Grothendieck,
written by Deligne, SGA~4\,$\frac12$, \emph{Cycle}, Sections~2.2.2 and~2.3.2.
The cycle-map compatibility with rational equivalence is also recorded
precisely in Milne, \emph{Lectures on \`Etale Cohomology}, version~2.21,
Section~23, Theorem~23.4 and the ensuing rational $\ell$-adic construction,
printed pp.~141--142.
The needed topological inputs are located in the proof below.
Neither A nor M is needed for these geometric and complexity statements.

\paragraph{Proof of rationality and generation of the Chow ring.}
On the nonempty principal open $f_1\ne0$ the defining equation solves for
\[
 y_1=\frac{t^2-t-\sum_{i=2}^h y_i f_i(x)}{f_1(x)}.
\]
Consequently this open is $D(f_1)\times\mathbb A^h_k$, with coordinates
$(x,t,y_2,\ldots,y_h)$. This proves rationality over every field, including
fields over which this particular open has no rational points. Smoothness,
geometric integrality and the smooth projection
$\pi:X\to\mathbb A^v_k$ of relative dimension $h$ are NN-C027.

Set $B=\mathbb A^v_k\setminus Y_\Phi$, $D=\coprod_{a,e}D_{a,e}$ and
$U=X\setminus D=\pi^{-1}(B)$. The open sets $D(f_i)$ cover $B$.
On each, solve for $y_i$ as above: the resulting $U\to B$ is Zariski
locally a product with $\mathbb A^h$. No globally chosen right inverse
to the row $(f_i)$ is asserted or needed.

Homotopy invariance gives $\operatorname{CH}^j(\mathbb A^v_k)=0$ for
$j>0$. Localization then gives $\operatorname{CH}^j(B)=0$ for $j>0$.
Affine-bundle surjectivity, which suffices here without any injectivity
claim, gives $\operatorname{CH}^j(U)=0$ for $j>0$. Since each component
of $D$ is an affine $h$-space, its only nonzero Chow group by dimension
is $\operatorname{CH}_h=\mathbb Z$. Localization for $D\subset X$ now
implies
\[
 \operatorname{CH}^j(X)=0\quad(j>0,\ j\ne v),\qquad
 \bigoplus_{a,e}\mathbb Z[D_{a,e}]\twoheadrightarrow
                            \operatorname{CH}^v(X).
\]
This calculation includes $j>v$: the possible source is then a lower
dimensional Chow group of $D$, which is zero, rather than an omitted
localization term.

Flat pullback by $\pi$ carries $[a]\in\operatorname{CH}^v(\mathbb A^v_k)$
to $[D_{a,0}]+[D_{a,1}]$. Both components have multiplicity one since
$t(t-1)$ has distinct simple roots over every field. As $[a]=0$ on
affine $v$-space, we obtain
\begin{equation}\label{eq:S016-opposite}
                     [D_{a,0}]+[D_{a,1}]=0.
\end{equation}
It follows that the $N$ classes $u_a$ generate $\operatorname{CH}^v(X)$.
Geometric integrality gives $\operatorname{CH}^0(X)=\mathbb Z[X]$.
The remaining issue is that the displayed generators have no further
relations, including torsion relations.

\paragraph{Proof of independence and of the cycle-class assertion.}
First pass to an algebraic closure and take $\mathbb Q_\ell$ coefficients.
The ordinary geometric cohomology of $X$, by the compact-support
calculation in NN-C027 and Poincar\'e duality in dimension $v+h$, is
\[
 H^0(X)=\mathbb Q_\ell,\qquad
 H^{2v}(X)=\mathbb Q_\ell(-v)^N,
 \qquad H^i(X)=0\ (i\ne0,2v).
\]
The relative compact-support splitting in NN-C027 is a sheaf identity
obtained from trace and involution projectors. Its geometric base-change
calculation works over every algebraically closed field, including
characteristic two; it does not infer these groups from Frobenius counts.
For $N=0$ the indicated second group is zero. Purity for
$Y_\Phi\subset\mathbb A^v$ and localization give
$H^{2v}(B)=0$: the relevant term to the right is
$H^{2v+1}_{Y_\Phi}(\mathbb A^v)=H^1(Y_\Phi,\mathbb Q_\ell(-v))=0$,
and affine space has no positive ordinary cohomology. On every
trivialization $D(f_i)$, affine-space invariance identifies
$\mathbb Q_\ell$ with $R(\pi|_U)_*\mathbb Q_\ell$.
The identification is the adjunction unit and hence glues; therefore
$H^*(U)=H^*(B)$ and in particular $H^{2v}(U)=0$.

Purity for the smooth closed pair $D\subset X$ and the cohomology
sequence with supports yield a surjection
\[
 \mathbb Q_\ell(-v)^{2N}=H^0(D,\mathbb Q_\ell(-v))
    =H_D^{2v}(X,\mathbb Q_\ell)
       \longrightarrow H^{2v}(X,\mathbb Q_\ell).
\]
The standard basis maps to the cycle classes of $D_{a,e}$, by the
fundamental-class normalization. Equation~\eqref{eq:S016-opposite}
also holds after applying the cycle map. Thus the $N$ images of
$u_a$ alone span the $N$-dimensional space $H^{2v}(X)$ and form a basis.
If an integral relation between the $u_a$ held over the original field
$k$, it would hold after base change and after applying this cycle map.
Its integer coefficients must therefore all be zero. This proves
integral independence, not only independence modulo torsion.

The presentation~\eqref{eq:S016-chow} now follows: a product of two
positive codimension classes lies in $\operatorname{CH}^{2v}(X)=0$
because $v\ge1$. The same bases prove~\eqref{eq:S016-cycle}, and all
other degrees are zero or the unit degree. Finally, an arbitrary
class $\sum_a n_a u_a$ is represented by the effective cycle
\[
 \sum_{n_a\ge0}n_a D_{a,0}
          +\sum_{n_a<0}(-n_a)D_{a,1}.
\]
Thus the monoid of effective codimension-$v$ cycle classes is the
entire group. This last fact uses affineness and rational equivalence;
it must not be treated as a pointed effective cone of a projective variety.

% Addition to root-owned NN-S016; not a separate registered finding.
% Standard amsmath/amssymb/amsthm environments only.
\begin{proposition}[NN-S016: complex homotopy type and the affine-space case]
Retain the sparse quadratic equations $f_1,\ldots,f_h$ of NN-C027,
with $v\geq1$ and
\[
 Y=V(f_1,\ldots,f_h)\simeq\coprod_N\operatorname{Spec}\mathbb Z,
 \qquad
 X=V\left(t^2-t-\sum_{i=1}^h y_i f_i(x)\right).
\]
The complex realization has homotopy type
\[
                      X(\mathbb C)\simeq\bigvee_{a=1}^{N}S^{2v},
\]
where the empty wedge means a point. In particular its integral
cohomology is $\mathbb Z$ in degree zero, $\mathbb Z^N$ in degree $2v$,
and zero otherwise, with every product of positive-degree classes zero.
Moreover,
\[
 N=0\quad\Longleftrightarrow\quad
 X\simeq\mathbb A^{v+h}_{\mathbb Z}
 \quad\Longleftrightarrow\quad
 X_{\mathbb C}\simeq\mathbb A^{v+h}_{\mathbb C}
 \quad\Longleftrightarrow\quad X(\mathbb C)\text{ is contractible}.
\]
The first isomorphism, when it exists, is proved by an explicit
polynomial formula from any identity $\sum_i a_i f_i=1$.
No Quillen--Suslin theorem or effective bound for such an identity is
required.
\end{proposition}

\begin{proof}
Write $\pi:X(\mathbb C)\to\mathbb C^v$ for the projection, and identify
$Y(\mathbb C)$ with its $N$ distinct points. First suppose $N>0$.
We prove the local models and the attaching maps before determining
the homotopy type.

Put $U=\mathbb C^v\setminus Y(\mathbb C)$ and
$A=\pi^{-1}(U)$. For $x\in U$ define the continuous column vector
\[
 a(x)=\frac{(\overline{f_1(x)},\ldots,\overline{f_h(x)})}
                  {\sum_i|f_i(x)|^2},\qquad f(x)\cdot a(x)=1.
\]
Every point of $A$ can uniquely be written
\[
       (x,t,y),\qquad y=a(x)t(t-1)+w,\quad f(x)\cdot w=0.
\]
Thus $A$ is the total space of the kernel vector bundle over
$U\times\mathbb C_t$. More explicitly, with $0\leq s\leq1$, set
\[
 t_s=(1-s)t,\qquad w_s=(1-s)w,\qquad
 y_s=a(x)t_s(t_s-1)+w_s.
\]
This is a strong deformation retraction onto the section $t=0,y=0$,
which is $U$. It is compatible with restriction to every open subset
of $U$. In particular projection identifies all intersection maps
below with the corresponding base-space inclusions. This is a
continuous construction; conjugation is not asserted to be algebraic.

For each $a\in Y(\mathbb C)$, reducedness says that the germs of the
$f_i$ generate the maximal ideal in the regular analytic local ring
of $\mathbb C^v$ at $a$. Their differentials therefore span the
cotangent space. Choose $v$ of them with independent differentials.
After reordering, the holomorphic inverse function theorem makes
\[
                    u_i=f_i(x),\qquad 1\leq i\leq v,
\]
coordinates on a neighborhood of $a$. In particular $h\geq v$.
Choose that neighborhood to map biholomorphically onto a small
Euclidean ball $B_a=\{\|u\|<\delta_a\}$ centered at zero, with these
neighborhoods pairwise disjoint. Each remaining holomorphic function
vanishes at zero and has an expression
\[
 f_j(x(u))=\sum_{i=1}^v b_{ji}(u)u_i\quad(j>v).
\]
For example the coefficients are obtained by integrating the partial
derivatives of $f_j(x(su))$ from $s=0$ to $s=1$, which is holomorphic
on the ball. The invertible triangular change of fiber coordinates
\[
 \eta_i=y_i+\sum_{j>v}b_{ji}(u)y_j\quad(1\leq i\leq v)
\]
turns the equation into
\[
                 \sum_{i=1}^v u_i\eta_i=t(t-1),
\]
with $h-v$ unused affine coordinates. Consequently
$V_a=\pi^{-1}(B_a)$ is analytically isomorphic to
\[
 \left\{(u,\eta,t):\sum_i u_i\eta_i=t(t-1),\ \|u\|<\delta_a\right\}
                   \times\mathbb C^{h-v}.
\]
Here and below $B_a$ also denotes its inverse image in the original
base coordinates.

The restriction $\|u\|<\delta_a$ does not change the homotopy type of
the full affine quadric $Q_{2v}:\sum u_i\eta_i=t(t-1)$.
Indeed take
\[
 c(u)=\min\left(1,\frac{\delta_a}{2\|u\|}\right),\quad c(0)=1,
 \qquad c_s(u)=1-s+s c(u).
\]
The maps $(u,\eta,t)\mapsto(c_s(u)u,c_s(u)^{-1}\eta,t)$ stay on the
quadric. At $s=1$ their image has $\|u\|<\delta_a$, and, if it starts
in that open subset, it remains there throughout the homotopy.
They supply a homotopy inverse to the inclusion of the restricted
quadric. This is a homotopy equivalence, not a claim that the whole
restricted subset is pointwise fixed during the homotopy.

The full quadric is linearly equivalent over $\mathbb C$ to
$\sum_{j=0}^{2v}z_j^2=1$: use $z_0=2t-1$ and, for each $i$,
$z_{2i-1}=\mathrm i(u_i+\eta_i)$, $z_{2i}=u_i-\eta_i$.
Writing $z=p+\mathrm i q$ with real vectors gives
\[
 |p|^2-|q|^2=1,\qquad p\cdot q=0.
\]
It is therefore diffeomorphic to the tangent bundle of $S^{2v}$,
by $(e,q)\mapsto\sqrt{1+|q|^2}\,e+\mathrm i q$ with
$|e|=1$ and $e\cdot q=0$. Scaling $q$ to zero retracts this bundle
onto its zero section. We have proved $V_a\simeq S^{2v}$, including
simple connectivity. On the intersection
\[
             W_a=A\cap V_a=\pi^{-1}(B_a\setminus\{a\}),
\]
the earlier projection retraction gives $W_a\simeq S^{2v-1}$.

Put $d=2v\geq2$. The small linking spheres give a basis of
$H_{d-1}(U;\mathbb Z)$, with no other reduced homology in $U$.
For precision, excise the complement of the balls from the pair
$(\mathbb C^v,U)$. Relative homology becomes the direct sum of the
pairs $(B_a,B_a\setminus\{a\})$. Each summand is $\mathbb Z$ in
degree $d$ and zero otherwise. As $\mathbb C^v$ is contractible,
the boundary map in the relative long exact sequence identifies
these $N$ generators with the linking-sphere basis in $H_{d-1}(U)$.
Naturality and the projection retractions prove that the inclusion
map
\[
     \bigoplus_a H_{d-1}(W_a;\mathbb Z)
                 \longrightarrow H_{d-1}(A;\mathbb Z)
\]
is an isomorphism. This verifies the attaching map on the needed
homology group; a dimension count alone would not suffice.

Apply Mayer--Vietoris to the open cover $X(\mathbb C)=A\cup B$,
where $B=\coprod_a V_a$ and $A\cap B=\coprod_a W_a$.
Its degree-$d$ portion is
\[
 0\longrightarrow\bigoplus_a H_d(V_a;\mathbb Z)
    \longrightarrow H_d(X(\mathbb C);\mathbb Z)
    \longrightarrow\bigoplus_a H_{d-1}(W_a;\mathbb Z)
    \xrightarrow{\ \sim\ }H_{d-1}(A;\mathbb Z).
\]
Thus the local quadric sphere classes give a basis of
$H_d(X(\mathbb C);\mathbb Z)=\mathbb Z^N$.
All other reduced homology groups vanish. At the bottom of the
sequence, including when $d=2$, the map
$H_0(A\cap B)=\mathbb Z^N\to H_0(A)\oplus H_0(B)=\mathbb Z\oplus\mathbb Z^N$
is $e_a\mapsto(1,-e_a)$ and is injective. This accounts for the
potential degree-one boundary term and proves connectivity.

Simple connectivity also requires a check. Add the $V_a$ one at a
time to $A$. They are simply connected and intersect the preceding
union in the path-connected $W_a$, since the $B_a$ are disjoint.
Van Kampen consequently expresses the fundamental group as
$\pi_1(A)$ modulo the normal subgroup generated by the images of
all $\pi_1(W_a)$, with paths to a common basepoint understood.
Projection identifies this quotient with the quotient obtained by
adding the contractible $B_a$ to $U$ in the base. The latter union
is $\mathbb C^v$, whose fundamental group is trivial. Thus
$X(\mathbb C)$ is simply connected. This argument treats $v=1$
without replacing its nontrivial puncture loops by a homology-only
argument.

The spaces involved have CW homotopy type. One may see this directly
from this open cover: $U$ deformation retracts onto a large
rectangular box with the interiors of $N$ small rectangular boxes
removed, by pushing radially away from each puncture and inward
from infinity. Subdividing a finite rectangular grid triangulates
that region. The $V_a$ and $W_a$ have the sphere homotopy types just
proved. The double mapping cylinder for the two-set open cover is
homotopy equivalent to $X(\mathbb C)$; replacing its vertex spaces
by these CW models and its attaching maps by cellular
approximations yields a CW model. Homotopy invariance of this gluing
and the open-cover equivalence are Hatcher, Propositions4G.1--4G.2,
pp.458--459; cellular approximation is Theorem4.8, p.349.

Since the space is simply connected and its homology below degree
$d$ vanishes, successive applications of Hurewicz give
$\pi_i(X(\mathbb C))=0$ for $i<d$ and
$\pi_d(X(\mathbb C))\simeq H_d(X(\mathbb C);\mathbb Z)$.
Represent the displayed $N$ homology basis elements by based sphere
maps and form their wedge. This map from $\bigvee_N S^d$ induces an
isomorphism on all integral homology groups. Both spaces are simply
connected and have CW homotopy type, so the homology form of
Whitehead's theorem makes the map a homotopy equivalence.
The cohomology-ring assertion follows: a product of positive-degree
classes would lie in degree $4v$, whose group vanishes.

If $N=0$, then $U=\mathbb C^v$ and the first explicit deformation
retraction already makes $X(\mathbb C)$ contractible. There is in
fact the stronger arithmetic isomorphism claimed in the statement.
The integral Boolean quotient in NN-C027 is the product of $N$
copies of $\mathbb Z$. If $N=0$ it is the zero ring, so the ideal
$(f_1,\ldots,f_h)$ equals the whole of $R=\mathbb Z[x]$.
Choose a polynomial column $a=(a_1,\ldots,a_h)$ with $f\cdot a=1$.
For a point $(x,t,y)$ of $X$, put
\[
                         z=y+a(2t-t^2).
\]
Then $f\cdot z=t^2-t+2t-t^2=t$. Conversely, from arbitrary
$(x,z)\in\mathbb A^{v+h}_{\mathbb Z}$ set
\[
             t=f(x)\cdot z,\qquad y=z+a(x)(t^2-2t).
\]
These polynomial formulas are inverse over $\mathbb Z$, and
$f\cdot y=t^2-t$ proves that the second lands in $X$.
Thus $X\simeq\mathbb A^{v+h}_{\mathbb Z}$.
If $N>0$, its nonzero degree-$2v$ integral cohomology forbids
contractibility and forbids an isomorphism with complex affine
space. Base change then also forbids the integral isomorphism.
This proves all the equivalences.
\end{proof}

\paragraph{Exact topological source interfaces.}
Allen Hatcher, \emph{Algebraic Topology}, author text
\url{https://pi.math.cornell.edu/~hatcher/AT/AT.pdf}:
Theorem1.20, p.43 (van Kampen);
Theorem2.20, p.119 (excision);
Section2.2, ``Mayer--Vietoris Sequences'', pp.149--150 (unnumbered
exact sequence and its inclusion maps);
Theorem4.32, pp.366--367 (Hurewicz);
Corollary4.33, pp.367--368 (homology Whitehead);
Theorem4.8, p.349 and Propositions4G.1--4G.2, pp.458--459
(CW replacement and open-cover gluing).
The cited statements and the local proof passages used here were
read in the author PDF. The explicit affine-quadric retraction,
local analytic reduction, linking basis and arithmetic polynomial
isomorphism are proved above.

\begin{proposition}[NN-S016: linear-degree integral coordinates in the empty case]
For the Boolean and verifier equations of NN-C027, suppose $N=0$.
There are integral polynomials $a_i(x)$ satisfying
\[
                      \sum_{i=1}^h a_i(x)f_i(x)=1,
 \qquad \deg a_i\leq v.
\]
Their coefficients have $O(v)$ bits. In the polynomial isomorphism
of the preceding proposition, the forward coordinate polynomials have
degree at most $v+2$ and the inverse coordinate polynomials have
degree at most $v+6$, with coefficient bit lengths $O(v+\log(h+1))$.
These are degree and height bounds; a fully expanded coordinate
list can still have exponentially many monomials.
\end{proposition}

\begin{proof}
Every equation has degree at most two and coefficient $\ell^1$-norm
at most three. All $v$ Boolean equations $B_j=x_j^2-x_j$ are retained.
Every non-Boolean verifier equation takes a value in $\{0,1,-1\}$
on a bit assignment. Since there are no satisfying assignments,
for each $b\in\{0,1\}^v$ choose an index $i(b)$ with
$f_{i(b)}(b)\in\{1,-1\}$. Put
\[
 e_b(x)=\prod_{j:b_j=1}x_j\prod_{j:b_j=0}(1-x_j),
 \qquad
 a_i^{(0)}(x)=\sum_{b:i(b)=i}f_i(b)e_b(x).
\]
Then $\deg a_i^{(0)}\leq v$, and evaluation on the Boolean cube gives
$\sum_i a_i^{(0)}(b)f_i(b)=1$. In particular
\[
 P(x)=\sum_i a_i^{(0)}(x)f_i(x)-1
\]
has degree at most $v+2$ and has zero multilinear remainder modulo
the $B_j$. This last assertion holds over $\mathbb Z$: the
multilinear polynomial with prescribed values on the Boolean cube
is uniquely $\sum_b P(b)e_b$, so all its values being zero implies
that it is the zero polynomial.

Here is explicit integral division with controlled degree and
height. For an exponent vector $\alpha$, write
$\epsilon_j=\min(\alpha_j,1)$. Telescope
\[
 x^\alpha-x^\epsilon=
 \sum_j\left(\prod_{l<j}x_l^{\epsilon_l}\right)
       (x_j^{\alpha_j}-x_j^{\epsilon_j})
       \left(\prod_{l>j}x_l^{\alpha_l}\right).
\]
The middle factor is zero for $\alpha_j\leq1$, and otherwise it is
\[
        (x_j^2-x_j)\sum_{r=0}^{\alpha_j-2}x_j^r.
\]
Apply this formula to every monomial of $P$. Since its multilinear
remainder is zero, the result is
\[
                         P=\sum_j q_j B_j,
 \qquad \deg q_j\leq v.
\]
Subtract $q_j$ from the multiplier of the corresponding retained
Boolean equation. The resulting $a_i$ have the asserted degree
and satisfy the ordinary polynomial identity $\sum_i a_i f_i=1$.
This is an identity in the unrestricted integral polynomial ring,
not just in its Boolean quotient.

For a polynomial $Q$, write $\|Q\|_1$ for the sum of the absolute
values of its coefficients. Since $\|e_b\|_1=2^{|\{j:b_j=0\}|}$,
\[
 \sum_i\|a_i^{(0)}\|_1\leq\sum_b\|e_b\|_1=3^v,
 \qquad \|P\|_1\leq3^{v+1}+1.
\]
The division formula introduces at most $|\alpha|\leq v+2$
quotient monomials per input monomial, counted with multiplicity.
Hence
\[
 \sum_j\|q_j\|_1\leq(v+2)(3^{v+1}+1),\qquad
 \|a_i\|_1\leq B_v:=3^v+(v+2)(3^{v+1}+1).
\]
Thus every integer coefficient has $O(v)$ bits.

For the forward isomorphism $z_i=y_i+a_i(2t-t^2)$, the degree is
at most $v+2$ and its coefficient norm is at most $1+3B_v$.
For the inverse, $t=\sum_i f_i(x)z_i$ has degree at most three and
coefficient norm at most $3h$. Therefore
\[
 y_i=z_i+a_i\left(\left(\sum_j f_jz_j\right)^2
                           -2\sum_j f_jz_j\right)
\]
has degree at most $v+6$ and coefficient norm at most
$1+B_v(9h^2+6h)$. These give the claimed bit bounds.
A degree-$v$ polynomial in $v$ variables can have
$\binom{2v}{v}\leq4^v$ monomials; neither the degree nor the
coefficient-height estimate bounds their count polynomially.
\end{proof}

\paragraph{Proof-system and geometric interpretation.}
The multipliers form an integral Nullstellensatz refutation of the
quadratic verifier equations with their Boolean axioms. They also
supply low-degree polynomial coordinates for the same smooth cubic
model's arithmetic affine-space isomorphism. This refines the
Boolean-quotient identity of NN-C009 by explicitly constructing the
ordinary ideal multipliers in this special empty Boolean-scheme
case. No polynomial-size expanded certificate, short arithmetic
circuit, or efficient isomorphism-finding algorithm follows from
this degree calculation alone. NN-C029 separately analyzes what E's
symbolic partition routine supplies for circuit descriptions.

\paragraph{Effectivity and complexity scope.}
On the encoded family of NN-C027, $N$ is the satisfying-assignment
count. The rank of the single positive-degree integral cohomology
group, and the number of spheres in this homotopy type, are therefore
counting-complete quantities; their parity is parity-counting-complete.
Complex contractibility, isomorphism to complex affine space, and
isomorphism to integral affine space are each coNP-complete on this
encoded family, because each is equivalent to unsatisfiability.
Their complementary questions are NP-complete. Every scheme under
discussion is affine; the distinction is whether it is isomorphic
to affine space. The polynomial identity $f\cdot a=1$ and the
resulting isomorphism need not have a polynomial-size description,
so this construction supplies no short coNP certificate. E evaluates
these decisions through the original Boolean formula, while A and M
are not needed for the homotopy or isomorphism statements.
No motivic or $\mathbb A^1$-homotopy equivalence is asserted by this
complex-topological proof.


\paragraph{Every field fiber, and stabilization.}
The integral isomorphism when $N=0$ gives an affine-space isomorphism
after every base change. If $N>0$, equation~\eqref{eq:S016-cycle} gives
nonzero $H^{2v}(X_{\bar k},\mathbb Q_\ell)$ over every field. This
forbids an affine-space isomorphism over that field or over its algebraic
closure. Multiplication by any affine space preserves this cohomology,
so it also forbids stable isomorphism to affine space. Thus the same
unsatisfiability equivalence holds for affine-space recognition over any
fixed field, and for stable affine-space recognition on the encoded family.

\paragraph{Complexity and GCT consequences.}
The algebraic construction has polynomial description length and retains
the original satisfying-assignment count $N$. The explicitly identified
Chow rank, reduced Betti number and sphere count are exactly $N$.
Counting accepting witnesses proves membership in \#P on the specified
encoding, and the construction is parsimonious, proving \#P-hardness.
Equation~\eqref{eq:S016-wedge} makes noncontractibility equivalent to
$N>0$ and contractibility equivalent to $N=0$. These give the NP and
coNP completeness assertions directly, without a general algorithm
for recognizing contractible varieties.

E computes $N$, and hence the compact homotopy description $(2v,N)$,
within its stated original-variable counting bound for the source CNF.
This does not output $N$ individual spheres or cycle generators in that
time when $N$ is exponentially large. For arbitrary verification circuits,
the original circuit-and-witness size replaces the CNF variable parameter
exactly as in NN-C028.

The example strengthens the geometric limit on what purity alone can
compute: the varieties are explicitly rational and simply connected,
all their cohomology is generated by displayed algebraic cycle types,
and the cycle-class map is an isomorphism. The rank still recovers \#SAT.
No unsolved assertion that Tate classes are algebraic is needed in this
family. A polynomial procedure extracting that rank would give
$\mathrm{FP}=\#\mathrm P$ and the familiar Toda collapse; an efficient
noncontractibility test would give $\mathrm P=\mathrm{NP}$. These remain
conditional algorithmic implications, with no unconditional lower bound.
Under A, the zeta factors from NN-C027 have their prescribed shifted
critical lines, while $N$ remains their lower multiplicity. M is unused.
The family supplies an actual cycle realization but does not identify
these cycles with a canonical GCT basis or the multiplicities of the
permanent and determinant orbit closures.

\paragraph{Geometric theorem and Grothendieck dictionary.}
The morphism $X_\Phi\to\operatorname{Spec}\mathbb Z$ is a smooth
geometrically integral rational sparse cubic model. For each field fiber,
the cycle classes of its two affine components above each Boolean section
satisfy exactly one relation, $[D_{a,0}]+[D_{a,1}]=0$, and no others.
Its integral Chow ring is the square-zero extension of $\mathbb Z$ by
the free rank-$N$ group in codimension $v$, and its rational $\ell$-adic
cycle map is an isomorphism. The positive-codimension effective class
monoid equals the whole Chow group; it is not an ordered representation
semiring with a pointed cone. On the complex fiber the corresponding
integral homotopy type is the wedge of $N$ even spheres. The proof is
the localization-and-purity calculation and the explicit local
quadric argument above, rather than an inference of topology from
point counts alone.

This places the statement in Grothendieck's cohomology with supports and
Poincar\'e duality, SGA~4 XVI~3.7/3.10(a), XVIII~3.2.5/3.2.6.2, and
the cycle class of SGA~4\,$\frac12$, \emph{Cycle}~2.2.2/2.3.2.
The primary SGA~4 re-edition is
\url{https://www.normalesup.org/~forgogozo/SGA4/16/16.pdf}, pp.~11--13,
and \url{https://www.normalesup.org/~forgogozo/SGA4/18/18.pdf},
pp.~77--78 (original pp.~585--586 for duality).
In particular the purity twist is $(-v)$; the minus sign is explicit
in the primary re-edition's Remark~3.10(a), p.~13.
The localization dictionary for cycles is Stacks Tags~02RX, 02TS and~02TY.
The completeness assertions refer to succinct input equations and basis
rank extraction, while the individual realized cycle classes have the
explicit geometric descriptions just given.
